\documentclass[11pt]{article}

\usepackage[a4paper,margin=1in]{geometry}
\usepackage{hyperref}
\renewcommand{\UrlFont}{\footnotesize}
\usepackage{enumitem}
\usepackage{titlesec}

\titleformat{\section}{\large\bfseries}{\thesection}{1em}{}
\setlist[itemize]{noitemsep, topsep=4pt}
\setlist[enumerate]{noitemsep, topsep=4pt}

\title{KCPC-CS120 \\ Problem Sheet \\ Combinatorics I}
\author{}
\date{}

\begin{document}

\maketitle

\noindent
Developed by the KCPC Competitive Programmes Department, this problem sheet accompanies the \textit{Combinatorics I} workshop.  
Problems are divided into three categories:
\begin{itemize}
    \item Integrated examples (discussed live)
    \item Core practice set (to be attempted during the workshop)
    \item Extended homework
\end{itemize}

\hrule
\vspace{1em}

% ----------------------------------------------------
\section*{I. Integrated Into Slides (Live Walkthrough)}

\begin{enumerate}
    \item \textbf{Project Euler 15 - Lattice Paths}  

    \vspace{0.3em}
    {\url{https://projecteuler.net/problem=15}}
                \vspace{0.8em}

    \item \textbf{LeetCode 62 - Unique Paths}  


    \vspace{0.3em}
    {\url{https://leetcode.com/problems/unique-paths/description/}}

              \vspace{0.8em}

    \item \textbf{Project Euler 53 - Combinatoric Selections}  
    
    \vspace{0.3em}
    {\url{https://projecteuler.net/problem=53}}
    
            \vspace{0.8em}

    \item \textbf{Project Euler 24 - Lexicographic Permutations}  
    
    \vspace{0.3em}
    {\url{https://projecteuler.net/problem=24}}

\end{enumerate}

\vspace{1em}
\hrule
\vspace{1em}

% ----------------------------------------------------
\section*{II. Core Workshop Practice Set}

These problems should be attempted during the workshop.  
They require careful modelling and correct algorithm selection.

For any questions that were not completed during the workshop, please attempt them independently afterwards.

\begin{enumerate}
    \item \textbf{CodeForces 1743A - Password}
    
    \vspace{0.3em}
    {\url{https://codeforces.com/problemset/problem/1743/A}}
                \vspace{0.8em}

    \item \textbf{CodeForces 1821A - Matching}
    
    \vspace{0.3em}
    {\url{https://codeforces.com/problemset/problem/1821/A}}
                \vspace{0.8em} 
          

    \item \textbf{Project Euler 31 - Coin Sums}  
    
    \vspace{0.3em}
    {\url{https://projecteuler.net/problem=31}}


    \vspace{0.8em}              

    \item \textbf{Codeforces C - Good Subarrays}  
    
    \vspace{0.3em}
    {\url{https://codeforces.com/problemset/problem/1398/C}}

    \vspace{0.8em}              

    \item \textbf{AtCoder D - Knight}  
    Very good! It is interesting, will require some knowledge of number theory
    \vspace{0.3em}
    {\url{https://atcoder.jp/contests/abc145/tasks/abc145_d}}
\end{enumerate}

\vspace{1em}
\hrule
\vspace{1em}

% ----------------------------------------------------
\section*{III. Extended Homework}

\textbf{Important.} The following problems are genuinely extremely difficult.  
They should only be attempted after completing all previously listed problems and assigned supplementary reading.  
They are intended as stretch challenges for those who wish to explore advanced modelling and optimisation techniques.

\medskip

\textbf{Yersin's Note.} 

I hope you enjoy working through these problems.

If you have any questions, please feel free to contact me or the Director of Competitive Programmes (Marcell).

\vspace{0.8em}

\vspace{1em}

\hrule
\vspace{1em}

\begin{enumerate}
    \item \textbf{Project Euler 483 - Repeated Permutation}
    
    \vspace{0.3em}
    {\url{https://projecteuler.net/problem=483}}
                \vspace{0.8em}

    \item \textbf{Leetcode 22 - Generate Parantheses}
    
    \vspace{0.3em}
    {\url{https://leetcode.com/problems/generate-parentheses/description/}}
                \vspace{0.8em} 
          

    \item \textbf{AtCoder T - Permutation}
    
    \vspace{0.3em}
    {\url{https://atcoder.jp/contests/dp/tasks/dp_t}}
\end{enumerate}

\noindent
\textbf{Guidance:}
\begin{itemize}
    \item Identify whether order matters: determine whether to use permutations or combinations.
    \item When direct counting overcounts, consider inclusion-exclusion or division by symmetry.
    \item For hard problems, try to find a bijection to a set that is easier to count.
    \item Focus on thinking of a nice argument before commiting anything to code.
\end{itemize}

\end{document}
